\section{Cognition: What the Flourishing of a Shared Life Consists In}
\label{p11:sec:cognition}

\noindent
The cycle is entered at the moment of cognition, the apprehension of what the flourishing of this shared life consists in. The task of this moment is not the measurement of a flourishing already defined, which belongs to evaluation, but the prior apprehension of what is to count as flourishing for these two, in this life: the framing within which all subsequent practice and assessment proceed. Three frameworks are drawn upon at this moment, the normative, the virtue-theoretic, and the phenomenological, and each contributes a distinct component to the cognition; the moment closes by stating the normative specification of a just flourishing on which the paper's later evaluation will rest.

\subsection{The normative component: flourishing as the unfolding of each}
\label{p11:sec:cog-normative}

The normative framework supplies the cognition that flourishing is not the satisfaction reported nor the conditions possessed but the substantive freedom of each party to unfold what that party has reason to value. The capabilities approach is the developed form of this cognition: the measure of a life is the range of beings and doings genuinely open to a person, the capabilities, and not merely the functionings achieved or the resources held \parencite{sen1999,nussbaum2011}. Transposed to the shared life, the cognition is that the flourishing of the relation is the substantive freedom of each party to unfold that party's own dynamics, the development of each into what each has it in them to become, sustained rather than foreclosed by the shared life.

This cognition fixes two things the other components and the later evaluation depend on. It fixes that flourishing is plural and not aggregate: it is the unfolding of each, and not a single quantity of relational satisfaction to which both contribute and in which their distinct developments are merged. A shared life has not flourished if it has produced a large sum of satisfaction unevenly distributed, with one party developed and the other foreclosed; the unfolding is of each, severally, or it is not the flourishing the cognition specifies. And it fixes that flourishing is a freedom and not a state: it is the openness of each party's development, the capability of unfolding, and not a particular condition attained, which preserves at the level of cognition the thesis that flourishing is not a state to be banked but a movement to be sustained.

The distinction the capabilities approach draws between a capability and a functioning bears directly on the cognition of a shared flourishing, and it is worth drawing out. A functioning is a being or doing achieved; a capability is the substantive freedom to achieve it, whether or not it is in fact achieved. The distinction matters because a shared life may produce achieved functionings, a party who is fed, housed, and occupied, while foreclosing the capability of which they are the husk, the party's substantive freedom to have done and become otherwise. A relation in which one party's functionings are provided for, but in which that party has no real freedom to pursue the developments they have reason to value, has supplied the functioning and foreclosed the capability, and it has not, on this cognition, sustained that party's flourishing. The point guards against a characteristic misreading of a shared life, in which the provision of a comfortable condition is taken for the flourishing of the one provided for; the cognition the normative framework supplies is that flourishing is the freedom to unfold, and that a comfort which forecloses the freedom is the mark of a foreclosure and not of a flourishing. The cognition also fixes what the later evaluation must therefore attempt, the assessment not merely of what each party has and does but of what each party is genuinely free to become within the shared life, which is the most demanding of the things the evaluation is asked to proxy and the one its proxies reach least adequately (\xref{p11:sec:ev-dynamics}).

\subsection{The virtue-theoretic component: flourishing as activity in accordance with cultivated excellence}
\label{p11:sec:cog-virtue}

The virtue-theoretic framework supplies a cognition the normative framework leaves implicit: that the unfolding of a party's dynamics is not a spontaneous expansion but an activity in accordance with excellences that are cultivated, and that the shared life is among the principal sites of their cultivation. In the Aristotelian account, eudaimonia is the activity of a life in accordance with excellence of character, and such excellence is acquired by habituation, the repeated performance of the corresponding actions until they settle into a disposition \parencite{aristotle2009}. In the Confucian account, the person is formed through sustained self-cultivation, much of it conducted in the ordering of the household and the practice of its daily observances \parencite{tu1985}. The conception of the person as constituted in relation, and of autonomy as itself developed rather than presupposed, is consonant with the self-determination tradition in psychology \parencite{ryan2000}. The cognition these traditions supply is that the flourishing of a shared life is not only the freedom of each to unfold but the activity of each, in accordance with excellences that the shared life cultivates, so that the relation is not merely the condition under which each develops but a principal medium of the development.

This component connects the cognition of flourishing to its practice, and so prepares the moment of practice (\xref{p11:sec:practice}). Because the excellences in accordance with which a life flourishes are cultivated by habituation, the cognition of what flourishing consists in already points toward the formation of habit and the ordering of a common life as the means of its pursuit. The virtue-theoretic cognition is, in this respect, the cognition that flourishing must be practised into being, and it forbids in advance any conception of the good life as a condition that could be secured otherwise than through the cultivation of disposition over time.

The virtue-theoretic component supplies, further, a cognition that the normative framework taken alone might miss: that in a shared life the excellences are cultivated jointly, and that some of them are excellences of the relation and not only of the parties severally. The patience, the attentiveness, the generosity, and the justice that a good shared life requires are dispositions formed in the parties by the practice of the shared life itself, each party habituated by the relation into the character the relation needs and, in being so habituated, developed; and there are, besides, excellences that belong to the pair, a way of attending to one another, a manner of repair after conflict, a shared bearing toward the world, which are dispositions of the relation, cultivated by its practice and possessed by neither party alone. The classical treatment of such friendship, as a relation in which each wills the good of the other for the other's sake and in which the friends are a condition of one another's flourishing, is Aristotle's \parencite{aristotle_ross}. The cognition that flourishing is activity in accordance with cultivated excellence thus carries, for a shared life, the recognition that the relation is a principal medium of the parties' development and the bearer of excellences of its own, and that to cognise the flourishing of a shared life is to cognise the excellences, of each party and of the pair, that the life is to cultivate. This recognition is what the moment of practice will take up, under the constraint that the cultivation of these excellences must guide and not legislate the relation's habit-forming dynamics (\xref{p11:sec:practice}).

\subsection{The phenomenological component: flourishing as lived and meant}
\label{p11:sec:cog-phenom}

The phenomenological framework supplies the cognition that flourishing is, before it is anything assessable, lived and understood: given as the felt quality and the apprehended meaning of a shared life, and not in the first instance as a fact about it. This component is a corrective to the other two as much as an addition. It holds that the normative specification of capabilities and the virtue-theoretic account of cultivated excellence, however just, remain abstract until they are referred to the lived texture of a particular shared life, in which alone the flourishing of these two has its concrete sense. The cognition of what flourishing consists in, for a given couple, is finally an interpretation of their shared life from within it, a hermeneutic apprehension of what their flourishing means, and not the application to them of a general specification from without.

The phenomenological component also supplies the cognition that what is apprehended at this moment is in large part unsymbolised. The felt sense of what a good shared life would be, for these two, precedes and exceeds its articulation; the cognition of flourishing is, at its root, the basal apprehension of \xref{p11:sec:appraisal}, conducted in a value not yet brought to words. This connects the moment of cognition to the account of appraisal and guards the moment against the supposition that to cognise flourishing is to formulate a definition of it. To cognise flourishing, at this moment, is in the first place to apprehend, in the lived and largely unsymbolised register, what the good of this shared life would be, and only secondarily to articulate it.

\subsection{The frameworks in concert and in tension}
\label{p11:sec:cog-concert}

The three components are drawn upon together, and they do not perfectly agree. The normative framework specifies flourishing as the unfolding of each, abstractly and for any couple; the virtue-theoretic framework locates it in the cultivation of excellence through shared habit; the phenomenological framework refers it to the lived and unsymbolised meaning of a particular life and resists its abstract specification. The tension is productive and is not to be resolved by the dominance of one component. The normative specification without the phenomenological reference is empty, a general account not yet referred to any actual life; the phenomenological apprehension without the normative specification is mute, a felt sense not yet articulated into anything that could guide practice or be assessed; and both without the virtue-theoretic component float free of the habituation through which alone a flourishing is brought into being. The cognition of flourishing is adequate when the three are held together, the normative specifying what is to be valued, the phenomenological referring it to the lived meaning of this life, and the virtue-theoretic binding it to the cultivation through which it is pursued.

\subsection{The normative specification of a just flourishing}
\label{p11:sec:cog-spec}

The moment of cognition closes by stating, for use in the paper's later evaluation, the specification of a just flourishing toward which the cognition tends. The specification gathers the cognition of this moment with the results of the prior paper.

\begin{claim}[The normative specification of a just flourishing]
A shared life flourishes justly when its cycle of value satisfies the following, jointly. The value generated in the shared life circulates and returns to those who generate it, no party being reduced to a source of value from which it is extracted without return. The unfolding of each party's own dynamics is sustained, neither party's development being foreclosed for the sake of the other's or of the relation's. And the cycle is sustained as a living and open movement, its value neither settled into a possessed condition nor cashed into a closed account, but generated anew in the turning. The first is the condition of justice, carried from the prior paper's account of generative justice and the return to the creator. The second is the eudaimonic condition, the unfolding of each. The third is the condition of the cycle's life, its sustenance as an open spiral rather than a closed repetition. A flourishing that meets the eudaimonic condition while failing the just is not a diminished flourishing but, as the prior paper established, no eudaimonia at all, the development of one purchased by the foreclosure of another. The three conditions are therefore not a list of desiderata to be traded against one another but the joint specification of a single thing, a just eudaimonia, and they furnish the normative basis against which the paper's later evaluation is conducted.
\end{claim}

\noindent
This specification is the product of the moment of cognition, and it is the standard the moment of evaluation will require. It is stated normatively, as what ought to be valued, and not yet as anything measured; the passage from this normative specification to an operational assessment, with all the concession that passage requires, is the work of \xref{p11:sec:evaluation}. The cycle now turns to the moment at which a flourishing so specified is pursued: the moment of practice.
